\section{Demo Interface}
\label{sec:demo-interface}

The demo begins like an ordinary streaming translator: the user selects a
language pair and starts a microphone or file stream. In AutoTerm mode there is
no glossary upload or domain selector. The system starts with a small working
set and updates it as translation context accumulates.

\paragraph{Walkthrough.}
The page first shows live partials while context accumulates. The working-set
badge then updates to the most similar slices, and retrieved pairs appear beside
the partial that used them. After a topic change, old and new slices can briefly
coexist until the sliding context reflects the new subject. The user changes no
control throughout this progression.

\begin{figure*}[t]
  \centering
  \includegraphics[width=0.72\textwidth]{ui_evidence_panel.png}
  \caption{The live En$\rightarrow$Zh interface after an automatic transition
  to medicine. Retrieved term pairs, source slices, and translation highlights
  expose the evidence used for each partial output.}
  \label{fig:ui}
\end{figure*}

\paragraph{Visible terminology evidence.}
The interface highlights retrieved terminology in the translation and lists
the corresponding source--target mappings. A compact status row shows the
active slices, term count, and retrieval latency. Thus a viewer can tell
whether a mistranslation came from missing retrieval evidence, an irrelevant
slice, or language-model decoding.

\paragraph{Extensibility.}
Operators register or replace topic slices offline; users see the benefit in
the next session without model training or UI changes. Fixed-glossary presets
and ad-hoc term overrides remain available for expert users, but they are not
needed for the zero-setup workflow.

Registration validates term pairs, builds an independent index, and adds one
manifest entry. Existing sessions retain their current snapshot, while new
sessions immediately see the updated candidate set. Thus an organization can
grow or revise its catalog without restarting the translation service or
rewriting the client application.

\paragraph{Availability.}
The MIT-licensed implementation, reviewer-friendly launcher, and screencast are
available in the
\href{https://github.com/luojiaxuan/autoterm-sst}{GitHub repository}. The
hosted entry page is
\url{https://luojiaxuan.github.io/autoterm-sst/}; live mode uses Qwen3-Omni
with vLLM continuous batching and the RASST retriever.
